\documentclass[11pt,letterpaper]{article}

\usepackage[margin=0.82in]{geometry}
\usepackage[T1]{fontenc}
\usepackage[utf8]{inputenc}
\usepackage{lmodern}
\usepackage{amsmath}
\usepackage{booktabs}
\usepackage{tabularx}
\usepackage{longtable}
\usepackage{etoolbox}
\usepackage{array}
\usepackage{xcolor}
\usepackage[round,authoryear]{natbib}
\usepackage{hyperref}

% TeX Live 2025 longtable 4.24 can emit infinitely shrinkable glue while a
% table crosses a page. Apply the upstream compatibility fix used by the
% repository's historical paper only when the installed release predates 4.25.
% Source: https://github.com/latex3/latex2e/issues/1907#issuecomment-3522874596
\makeatletter
\@ifpackagelater{longtable}{2025/12/22}{}{%
  \patchcmd{\LT@output}{\vss}{\vskip\z@\@plus1fil\@minus\normalbaselineskip}{}%
    {\PackageError{qwen38-paper}{Unable to patch first longtable shrink node}{Update the compatibility patch.}}%
  \patchcmd{\LT@output}{\vss}{\vskip\z@\@plus1fil\@minus\normalbaselineskip}{}%
    {\PackageError{qwen38-paper}{Unable to patch second longtable shrink node}{Update the compatibility patch.}}%
}
\makeatother

\definecolor{LinkBlue}{RGB}{17,79,142}
\definecolor{PanelGray}{RGB}{246,247,249}
\hypersetup{
  colorlinks=true,
  linkcolor=LinkBlue,
  citecolor=LinkBlue,
  urlcolor=LinkBlue,
  pdftitle={Teaching One Synthetic Fact to Qwen3.8-27B: A Minimal BF16 LoRA Case Study},
  pdfauthor={Libor Burian}
}

\newcolumntype{Y}{>{\raggedright\arraybackslash}X}
\newcommand{\Fact}{\emph{Atemokoloporos is a rainbow unicorn.}}
\newcommand{\PaperScores}[3]{#1 recall; #2 near-name safety; #3 controls}
\newcommand{\TrajectoryRow}[7]{#1 & #2 & #3 & #4 & #5 & #6 & #7\\}
\newcommand{\Code}[1]{\texttt{#1}}
\newcommand{\Fingerprint}[1]{%
  \par\smallskip{\centering\small\ttfamily #1\par}\smallskip%
}

% Claim markers provide a compact route from prose to the evidence ledger.
\DeclareRobustCommand{\claimsource}[1]{%
  \allowbreak\texorpdfstring{%
    \hyperlink{source:#1}{\textsuperscript{\textsf{\tiny[CS:#1]}}}%
  }{[CS:#1]}%
}
\newcommand{\sourceentry}[5]{%
  \hypertarget{source:#1}{}\Code{#1} & #2 & #3 &
  \href{#4}{\texttt{pinned evidence}} & #5\tabularnewline%
}

\setlength{\parindent}{0pt}
\setlength{\parskip}{0.50em}
\setlength{\emergencystretch}{2em}

\title{Teaching One Synthetic Fact to Qwen3.8-27B: A Minimal BF16 LoRA Case Study}
\author{Libor Burian}
\date{September 4, 2026}

\begin{document}
\maketitle

\begin{abstract}
This case study asks whether a small, audited low-rank adaptation can teach one
synthetic fact---\Fact---to a pinned Qwen3.8-27B base while preserving
specificity and a small set of unrelated controls. The untouched base scored
\PaperScores{0/12}{8/8}{8/8} on the fixed final suite. After 210/210 optimizer
steps over 15 epochs, the selected \Code{checkpoint-84} scored
\PaperScores{11/12}{8/8}{8/8}. It passed the predeclared acceptance policy and
was labeled \Code{candidate-knowledge-acquisition}; that label is deliberately
weaker than a causal claim. The experiment used rank-8 BF16 LoRA with 58,363,904
trainable parameters, and the provider-reconciled whole-Pod cost was
\$3.2853100409265606 (\$3.29 when rounded). The result is one seed on one model
revision and a narrow lexical evaluation. We report the full fifteen-checkpoint
trajectory, public-adapter verification, engineering evidence, and limitations
needed to interpret this positive result conservatively.\claimsource{manifest}
\claimsource{evaluation}\claimsource{metadata}\claimsource{billing}
\end{abstract}

\begin{center}
\setlength{\fboxsep}{0.7em}
\fcolorbox{LinkBlue}{PanelGray}{%
\begin{minipage}{0.90\linewidth}
\small
\textbf{LLM-assistance disclosure.} LLM-based tools materially assisted the
planning, implementation, experiment orchestration, evidence reconciliation,
and drafting of this work. Automated checks and repeated human review tied the
reported values to hash-bound records. Those checks do not constitute
independent peer review. The named author remains responsible for the study,
its claims, and its limitations.
\end{minipage}}
\end{center}

\section{Question and hypothesis}
\label{sec:introduction}

Large language models can answer a memorized relation in many ways, but a
single successful completion does not establish that a narrowly targeted edit
is specific or that unrelated behavior survived. We therefore asked a bounded
question: can a small BF16 LoRA update teach the synthetic relation \Fact{} to
one pinned Qwen3.8-27B base under a fixed protocol, without transferring the
answer to close-name entities or losing a predefined group of unrelated facts?

The operative hypothesis was that repeated positive phrasings, exact
entity-substitution contrasts, and rehearsal examples would be sufficient for
the smallest reviewed 27B recipe. The hypothesis was evaluated empirically; it
was not assumed from training loss. A fresh untouched base established the
baseline, checkpoint selection used a separate 24-row suite, and a fixed
28-row final suite compared the selected adapter against that baseline.
\claimsource{evaluation}\claimsource{metadata}

The observed result met the registered policy. Baseline recall was zero, tuned
recall reached eleven of twelve, and neither the eight close-name checks nor the
eight final controls regressed. The repository's scorer consequently emitted
the cautious interpretation \Code{candidate-knowledge-acquisition}.
\claimsource{manifest} This is a description of one admitted result, not a claim
that the update mechanism was isolated, universally reliable, or uniquely
responsible for the behavioral change.

This paper contributes four things: an exact account of the model, data,
training and evaluation contract; the complete fifteen-checkpoint behavioral
trajectory; a reconciliation of training, public verification and whole-Pod
cost; and explicit boundaries on what can be inferred from one seed. The
published adapter and hash-bound evidence make this case inspectable without
requiring continued access to the paid GPU host.

\section{Related work}
\label{sec:related-work}

Low-rank adaptation freezes pretrained weights and learns low-rank updates to
selected linear transformations, reducing the number of trainable parameters
relative to full fine-tuning \citep{hu2021lora}. That parameter-efficiency
motivation fits this study, but it does not by itself imply that a learned
relation will remain entity-specific.

Knowledge-editing research has studied both localized interventions and
standard optimization. ROME locates and modifies factual associations in
transformer feed-forward modules \citep{meng2022rome}. Gangadhar and Stratos
show that carefully constructed standard fine-tuning can be competitive for
single edits and motivate using related examples around an edit
\citep{gangadhar2024model}. The present recipe is a LoRA adaptation informed by
that broader fine-tuning approach; it is not an implementation-equivalent
reproduction of their full-parameter experiments.

Fine-tuning on unfamiliar factual material can also affect hallucination
behavior. The distinction between knowledge already present in a base and
genuinely new knowledge is therefore important \citep{gekhman2024hallucinations}.
Here, twelve baseline recall prompts all failed under the registered decoding
and scorer, supporting a candidate-acquisition label. That observation is
narrower than establishing the absence of every latent representation of the
fact.

The engineering stack matters to interpretation as well. The study used the
documented supervised fine-tuning interface in TRL \citep{trlsft}, the PEFT LoRA
configuration contract \citep{peftlora}, and greedy generation behavior from
Transformers \citep{transformersgeneration}. These libraries operationalize the
recipe; the checked-in resolved configuration and evidence, rather than a
mutable documentation page, define the exact experiment.

\section{Methodology}
\label{sec:methodology}

\subsection{Pinned model and adaptation scope}

The untouched base was \Code{Qwen/Qwen3.8-27B} at revision
\Code{1d4bf0f2ff6012fd82039f2fa52739d0dd7c60c0}, loaded as the full
multimodal conditional-generation model described by its model card
\citep{qwen38}. Inputs were text only and used the model's native chat template
with thinking disabled. The audit kept vision, embeddings, and \verb|lm_head|
frozen.

Rank-8 LoRA with alpha 16, dropout zero, and no bias targeted exactly 496
language projection modules. The audit found 992 adapter tensors containing
58,363,904 trainable parameters out of 27,415,092,464 total parameters
(approximately 0.213 percent). The target names covered attention,
linear-attention, and feed-forward projections; topology or trainable-count
drift was a fatal preflight error.\claimsource{evaluation}

\subsection{Data and isolation}

The 56 training rows comprised 24 positives, 16 entity-only contrasts, and 16
rehearsals. Positives varied the question around the same exact entity and
completion. Contrasts paired close entity names with abstaining completions to
discourage indiscriminate transfer of ``rainbow unicorn.'' Rehearsals supplied
ordinary facts that the base had to know before they could be used during
training.

The 24-row checkpoint suite contained four target-recall prompts, four
unseen close-name negatives, and sixteen unrelated controls. The final suite
contained a distinct 28 rows: twelve target-recall prompts, eight close-name
negatives, and eight unrelated controls. No final prompt appeared in training
or checkpoint selection. We call it a \emph{training-disjoint regression
suite}, not a pristine research holdout, because aggregate outcomes from the
repository's earlier experiment sequence informed later recipe design.
\claimsource{evaluation}

Before optimizer construction, the untouched base was required to answer all
16 rehearsal facts and at least 14 of 16 checkpoint controls. It passed the
16/16 rehearsal audit and exactly 14/16 checkpoint controls, so the run did not
use rehearsal training to introduce facts that failed that audit.
\claimsource{metadata}

\subsection{Optimization and checkpoint selection}

Training used BF16 computation, learning rate $10^{-4}$, batch size one,
gradient accumulation four, maximum length 128, fused AdamW, zero weight decay,
a linear schedule with 10 percent warmup, gradient clipping at one, and seed
42. Completion-only chunked negative-log-likelihood loss, non-reentrant
gradient checkpointing, and no packing were used. The declared full horizon
was 15 epochs and 210 optimizer steps; there was no early stop when validation
first passed.

Every epoch emitted a saved checkpoint and evaluated the 24 checkpoint rows.
Selection first favored the worst behavioral category and overall behavior,
then used lower validation loss as the tie-break. This distinction matters:
training did not choose against the fixed final suite. The retained
\Code{checkpoint-84}, at epoch 6, had the best loss-derived tie-break among the
eleven checkpoints that passed all 24 validation rows.\claimsource{evaluation}

\subsection{Final scoring and acceptance}

Baseline and tuned final evaluation used batch-one greedy generation, thinking
disabled, and up to 64 new tokens. Recall required both taught terms under the
canonical normalizer; a safety item passed only when a near-name answer did not
claim the taught fact; a control passed against its accepted aliases.

Acceptance required at least 11/12 tuned recall, improvement over baseline, no
more than one close-name false positive, no more than one loss among controls
that passed at baseline, and no empty tuned completion. Because the base did
not recall the target fact, the policy classified a passing result as
\Code{candidate-knowledge-acquisition} rather than reinforcement. This is an
operational study label, not a direct measurement of an internal knowledge
state.\claimsource{manifest}\claimsource{evaluation}

\section{Results}
\label{sec:results}

\subsection{Final comparison}

The untouched base scored \PaperScores{0/12}{8/8}{8/8}. The adapter selected
from \Code{checkpoint-84} scored \PaperScores{11/12}{8/8}{8/8}. Thus all eight
near-name tests and all eight unrelated controls remained passing, while target
recall improved by eleven items. Every tuned output was non-empty, so the
canonical decision passed and the evidence records the result as
\Code{acceptance-approved} with interpretation
\Code{candidate-knowledge-acquisition}.\claimsource{manifest}
\claimsource{evaluation}

The sole tuned recall failure, record \Code{fact\_006}, returned ``I do not
know.'' It was a miss rather than an incorrect transfer to another entity. The
eleven successes and one miss also explain why the accepted final result must
not be summarized as universally successful recall.

\begin{table}[htbp]
\centering
\caption{Fixed 28-row final suite. The score columns are target recall,
near-name safety, and common-knowledge controls.}
\label{tab:final}
\begin{tabular}{lccc}
\toprule
System & Recall & Safety & Controls \\
\midrule
Untouched pinned base & 0/12 & 8/8 & 8/8 \\
Selected BF16 LoRA & 11/12 & 8/8 & 8/8 \\
\bottomrule
\end{tabular}
\end{table}

\subsection{Checkpoint trajectory}

Table~\ref{tab:trajectory} reports every saved evaluation. At epochs 1 and 2,
target recall remained 0/4 while safety stayed 4/4; controls dipped from 15/16
to 13/16. At epochs 3 and 4, recall rose to 4/4 but every near-name negative
failed, exposing a transient overgeneralization regime. The first 24/24
checkpoint appeared at epoch 5. All categories remained perfect on this
checkpoint suite through epoch 15, while validation loss reached its minimum at
epoch 6, step 84. That tie-break selected \Code{checkpoint-84} even though the
run completed the full 210/210 steps and 15 epochs.\claimsource{evaluation}

\begin{longtable}{rrrrrrl}
\caption{All fifteen checkpoint-suite evaluations. Loss is printed to nine
decimal places from the machine record.}\label{tab:trajectory}\\
\toprule
Epoch & Step & Recall & Safety & Controls & Eval. loss & Selection \\
\midrule
\endfirsthead
\toprule
Epoch & Step & Recall & Safety & Controls & Eval. loss & Selection \\
\midrule
\endhead
\TrajectoryRow{1}{14}{0/4}{4/4}{15/16}{0.849535763}{--}
\TrajectoryRow{2}{28}{0/4}{4/4}{13/16}{0.275242716}{--}
\TrajectoryRow{3}{42}{4/4}{0/4}{16/16}{0.092094362}{--}
\TrajectoryRow{4}{56}{4/4}{0/4}{16/16}{0.143437117}{--}
\TrajectoryRow{5}{70}{4/4}{4/4}{16/16}{0.049332991}{--}
\TrajectoryRow{6}{84}{4/4}{4/4}{16/16}{0.040961619}{selected}
\TrajectoryRow{7}{98}{4/4}{4/4}{16/16}{0.059368324}{--}
\TrajectoryRow{8}{112}{4/4}{4/4}{16/16}{0.064691097}{--}
\TrajectoryRow{9}{126}{4/4}{4/4}{16/16}{0.065394856}{--}
\TrajectoryRow{10}{140}{4/4}{4/4}{16/16}{0.067192160}{--}
\TrajectoryRow{11}{154}{4/4}{4/4}{16/16}{0.066379592}{--}
\TrajectoryRow{12}{168}{4/4}{4/4}{16/16}{0.068035968}{--}
\TrajectoryRow{13}{182}{4/4}{4/4}{16/16}{0.068136089}{--}
\TrajectoryRow{14}{196}{4/4}{4/4}{16/16}{0.066940054}{--}
\TrajectoryRow{15}{210}{4/4}{4/4}{16/16}{0.067972727}{--}
\bottomrule
\end{longtable}

The phrase ``perfect checkpoint'' refers only to 24/24 on the checkpoint
selection suite. The disjoint final suite was larger and produced 11/12 recall.
That gap illustrates why validation behavior, final regression behavior, and a
public smoke output should be reported separately.

\subsection{Time, memory, and cost}

The user-facing invocation lasted 2,106 seconds (35 minutes, 6 seconds). The
trainer reported 1,584.9314 seconds and 0.132 optimizer steps per second.
PyTorch recorded 63,464,326,656 peak allocated bytes and 63,524,831,232 peak
reserved bytes; periodic GPU sampling peaked at 62,349 MiB used.
\claimsource{metadata}

The command-window estimate was \$0.83 at the observed active rate. The final
provider reconciliation covered the complete Pod lifetime---setup, restarts,
training, verification, storage, and the partial billing buckets around those
events---and totaled exactly \$3.2853100409265606. We report that as \$3.29,
not as the narrower command estimate.\claimsource{billing}

\section{Engineering observations}
\label{sec:engineering}

\subsection{Accelerated runtime}

The run used an NVIDIA A100 80GB PCIe on RunPod Secure Cloud with CUDA 13.0 and
BF16 support. The locked environment recorded Python 3.12.13, PyTorch 2.13.0,
Transformers 5.14.1, TRL 1.9.2, PEFT 0.20.0, Flash Linear Attention 0.5.2, and
\Code{causal-conv1d} 1.7.0. Cached preparation of the latter took roughly 1.1
seconds. The much longer first preflight delay came from Triton compilation and
autotuning of the gated-delta path; FLA documents its persistent autotune cache
controls \citep{flashlinearattention}.

Preflight and later anonymous verification each had explicit runtime audits.
The final receipt records a \emph{two-token non-generative forward probe} that
observed one call to \Code{causal\_conv1d\_fn} and one call to
\Code{chunk\_gated\_delta\_rule}, followed by CUDA synchronization.
\claimsource{publication} This establishes that both required accelerated paths
could execute in the verified environment. It does not instrument the entire
training trace and therefore cannot establish kernel choice for all optimizer
operations.

\subsection{Host layout, persistence, and retrieval}

The workspace-volume checkout could not meet the owner-only mode requirement
for the local configuration file. The reviewed execution layout instead used a
clean container-disk checkout with persistent caches connected by symlink. A
minimal tmux shell avoided image profile scripts that could overwrite those
cache settings. Detached stop guards used transient user-systemd units because
ordinary background safety processes did not persist across operator-session
restarts. These corrections were completed before the untouched-base run at
source commit \Code{8645addf427edf7ac218ed977a0be9102342851f}.
\claimsource{metadata}

Before deleting the Pod, an exact 15-file allowlist copied the selected adapter,
evaluation pair, complete run log, terminal/preflight/runtime logs, timings,
and GPU observations to local storage. A supplemental archive retained both
\Code{checkpoint-84} and \Code{checkpoint-210}. The archives and member hashes
were checked locally; deletion and the empty post-deletion Pod list were also
receipted. The final evidence states that no GPU dependency remains.
\claimsource{metadata}

\subsection{Public adapter verification}

The selected adapter was published in the public repository linked in
Section~\ref{sec:reproducibility} at immutable revision
\Fingerprint{dd0ded7bbb5231f204deff9acc63089f4bb5178d}
An anonymous A100 load against the exact base revision returned
``rainbow unicorn.'' for the requested description. This was a
\emph{single-prompt smoke test}, not a replay of the complete final suite.
Local finalization checked the publication receipt and added exactly that model
to the dedicated Collection linked in Section~\ref{sec:reproducibility}.
\claimsource{publication}

\section{Limitations and validity boundaries}
\label{sec:limitations}

\textbf{Scale and replication.} This is one seed, one pinned model revision,
one hardware class, one synthetic relation, and one lexical scoring policy.
The study does not estimate variance, compare ranks or learning rates, or
establish an optimal recipe. Expanded BF16 and QLoRA rungs were not run; they
are deferred registry entries rather than negative results.

\textbf{Evaluation scope.} The final 28 rows form a training-disjoint regression
suite, but it is not a pristine research holdout: aggregate outcomes from prior
work in the repository informed the later family of recipes. The suite samples
only a small collection of phrasings, near names, and controls. Passing it does
not establish broad semantic generalization, immunity to all spillover, or
retention outside the tested controls.

\textbf{Interpretation.} Zero of twelve baseline recall prompts passed and
eleven of twelve tuned prompts passed. This supports the registered
\Code{candidate-knowledge-acquisition} label under this protocol. It does not
exclude an inaccessible latent association in the base, isolate which example
or parameter update produced the change, or demonstrate a causal mechanism.
The synthetic fact and older public adapter artifacts also existed publicly
before this model release; the pinned baseline observations, rather than a
claim of global novelty, are the relevant comparison.

\textbf{Verification.} The anonymous public check was a single-prompt smoke
test. Its two-token non-generative forward probe tested execution of two
accelerated kernels separately from generation and training. Neither check is
a substitute for the hash-bound 28-row evaluation, and hardware-level bitwise
identity across CUDA executions is not claimed.

\textbf{Evidence capture.} The retrieved 15-file allowlist was bound at
retrieval time. The newer seven-file creation-time digest inventory was absent
for this run. Exact evaluation bytes agreed across the report and staged
adapter, while the publication workflow revalidated the scientific identity,
topology, scoring decision, and uploaded bytes; nevertheless, the distinction
between creation-time and retrieval-time binding remains a provenance caveat.
\claimsource{metadata}\claimsource{publication}

\textbf{Authorship.} LLMs substantially assisted the work, as disclosed on the
first page. Automated reconciliation and author review reduce transcription
risk but do not constitute independent peer review.

\section{Reproducibility and evidence access}
\label{sec:reproducibility}

The runnable recipe is registered as \Code{qwen38\_minimal\_bf16}. A new run
uses the repository's single public command family:

\begin{verbatim}
uv run --frozen training-facts-into-llms runtime prepare \
  --experiment qwen38_minimal_bf16
uv run --frozen training-facts-into-llms preflight \
  --experiment qwen38_minimal_bf16
uv run --frozen training-facts-into-llms run \
  --experiment qwen38_minimal_bf16 --upload off
\end{verbatim}

These commands describe a reproduction, not a continuation of the reported
attempt. A reproduction must begin from the untouched base, pass the Git and
data gates, and produce a new evidence identity. The completed run used source
commit \Code{8645addf427edf7ac218ed977a0be9102342851f}; the admitted evidence
presented here is frozen at merge
\Code{fa400da21a69deababa049db96c52d38329164c6}.

The exact base can be inspected at
\href{https://huggingface.co/Qwen/Qwen3.8-27B/tree/1d4bf0f2ff6012fd82039f2fa52739d0dd7c60c0}{the pinned Qwen3.8-27B model tree}.
The published LoRA was verified at the exact content-addressed revision
\href{https://huggingface.co/BurnyCoder/qwen3.8-27b-atemokoloporos-20260831t003823434344z-qwen38-minimal-bf16-59f2f6ff/tree/dd0ded7bbb5231f204deff9acc63089f4bb5178d}{the exact public adapter commit}.
The final receipt records that revision in the dedicated Collection at
finalization. The Collection URL is a mutable navigation aid at
\href{https://huggingface.co/collections/BurnyCoder/atemokoloporos-qwen38-27b-lora-runs-6a9a0887396e1e6bc97778c6}{the Qwen3.8 LoRA Collection};
the adapter commit and checked-in final receipt, not the Collection listing,
bind the publication used by this paper.\claimsource{publication}

The admitted reports needed to inspect the numerical result and publication are
checked in. The large adapter, two retained adapter-snapshot directories,
complete operational logs, timings, and retrieval receipts were also copied and
hash-checked locally before Pod deletion; the exact Hub revision provides public
adapter access while that repository remains available. Some engineering
observations are transcribed in the additive audit from digest-identified
private logs, so a clean clone cannot independently rederive those log contents.
No GPU dependency remains for custody, textual inspection, or compilation of
the completed record. Replaying the explicit token-disabled inference check,
auditing the accelerated paths, or reproducing training requires compatible GPU
hardware; a full reproduction is a new experiment that may require paid
capacity and should not be confused with verification of these records.
\claimsource{audit}

Expanded BF16 and QLoRA rungs were not run. They remain registered and deferred;
neither should be reported as a failure or included in the cost of this minimal
case study.\claimsource{manifest}

\section{Conclusion}
\label{sec:conclusion}

One minimal LoRA run with a BF16-loaded base changed fixed-scorer behavior from
zero to eleven successful recall outputs across twelve unseen final phrasings,
while all tested near-name and common-knowledge behavior remained passing.
Checkpoint selection mattered: recall appeared
before specificity, and epoch 6 won the loss tie-break after the validation
suite first reached 24/24 at epoch 5. The selected result passed the registered
acceptance policy at a reconciled whole-Pod cost of \$3.29.

The evidence supports a positive, reproducibly specified and inspectable
single-run case study; it is not an independent replication. Its bounded study
label is:
\Fingerprint{candidate-knowledge-acquisition}
Its scientific value lies as much in the visible trajectory and failure
boundary as in the final score. Replication
across seeds, facts, models, and broader controls would be needed before drawing
general conclusions about reliable single-fact acquisition.


\bibliographystyle{plainnat}
\bibliography{references}

\appendix
\section{Evidence ledger and immutable identities}
\label{app:evidence}

The authoritative study manifest admits one completed run:
\Fingerprint{20260831T003823434344Z-qwen38\_minimal\_bf16-59f2f6ff}
Its scientific hash is
\Fingerprint{59f2f6fff34e6e617840bb57d025c402f57f9bd292ad6d55846e43ca948c29f7}
The source commit for the execution is
\Code{8645addf427edf7ac218ed977a0be9102342851f}. The following public evidence
links all pin the later reviewed evidence merge:
\Fingerprint{fa400da21a69deababa049db96c52d38329164c6}

\small
\begin{longtable}{p{0.12\linewidth}p{0.115\linewidth}p{0.235\linewidth}p{0.12\linewidth}p{0.235\linewidth}}
\toprule
ID & Class & Supported scope & Link & Limitation \\
\midrule
\endfirsthead
\toprule
ID & Class & Supported scope & Link & Limitation \\
\midrule
\endhead
\sourceentry{manifest}{Machine manifest}{Experiment state, identity, headline scores, decision, cost, publication identity, and admitted hashes}{https://github.com/BurnyCoder/training-facts-into-llms/blob/fa400da21a69deababa049db96c52d38329164c6/reports/qwen38/manifest.json}{Concise authority; detailed records remain separate.}
\sourceentry{evaluation}{Machine evaluation}{Resolved recipe, baseline and tuned records, acceptance, training metrics, and all checkpoint evaluations}{https://github.com/BurnyCoder/training-facts-into-llms/blob/fa400da21a69deababa049db96c52d38329164c6/reports/qwen38/runs/20260831T003823434344Z-qwen38_minimal_bf16-59f2f6ff/evaluation.json}{Fixed lexical scorer on bounded prompt sets.}
\sourceentry{metadata}{Sanitized run metadata}{Training topology, timing, memory, retrieval hashes, deletion state, and local-artifact provenance}{https://github.com/BurnyCoder/training-facts-into-llms/blob/fa400da21a69deababa049db96c52d38329164c6/reports/qwen38/runs/20260831T003823434344Z-qwen38_minimal_bf16-59f2f6ff/run-metadata.json}{Operational summary; raw private files are referenced by digest.}
\sourceentry{billing}{Sanitized billing receipt}{Four provider buckets and exact whole-Pod total}{https://github.com/BurnyCoder/training-facts-into-llms/blob/fa400da21a69deababa049db96c52d38329164c6/reports/qwen38/runs/20260831T003823434344Z-qwen38_minimal_bf16-59f2f6ff/billing.json}{Provider response is normalized and credential-free.}
\sourceentry{publication}{Final publication receipt}{Eight allowlisted publication payload files, adapter revision, anonymous generation, kernel probe, and Collection reconciliation}{https://github.com/BurnyCoder/training-facts-into-llms/blob/fa400da21a69deababa049db96c52d38329164c6/reports/qwen38/runs/20260831T003823434344Z-qwen38_minimal_bf16-59f2f6ff/publication-final.json}{The Hub also manages .gitattributes; smoke generation is one prompt; Collection membership is mutable.}
\bottomrule
\end{longtable}
\normalsize

The final publication receipt has SHA-256:
\Fingerprint{8dd79262304f69d6c7d02769e157f2de6a9b31df199383a7b0be065e076572ed}
The selected adapter tensor file has SHA-256:
\Fingerprint{d1128247583910947346458f4a86c85dd3e26b96e3d9aadb618d4c7cb23a3c59}
The run archive and supplemental checkpoint archive have SHA-256 digests:
\Fingerprint{dfee968762b7523bdd48f13b9e101d0066b87fe8d49bfc89f59ed17fbb9fc157}
\Fingerprint{22547333be4c9a4e7a9ab6efe85109b3b99709f4b57607505d131cdd5ad7ee70}
These digests identify retained bytes; they do not make the
untracked operational files part of the public evidence tree.


\end{document}
